\documentclass[11pt]{article}
\usepackage[T1]{fontenc}
\usepackage{lmodern}
\usepackage{microtype}
\usepackage{amsmath}
\usepackage{amssymb}
\usepackage{upgreek}
\usepackage{booktabs}
\usepackage{xcolor}
\usepackage[margin=1in]{geometry}
\setlength{\parindent}{0pt}
\setlength{\parskip}{0.65em}
\setlength{\jot}{0.45em}
\emergencystretch=2em
\allowdisplaybreaks[4]
\DeclareUnicodeCharacter{0393}{\ensuremath{\Gamma}}
\DeclareUnicodeCharacter{0391}{A}
\DeclareUnicodeCharacter{0392}{B}
\DeclareUnicodeCharacter{0394}{\ensuremath{\Delta}}
\DeclareUnicodeCharacter{0395}{E}
\DeclareUnicodeCharacter{0396}{Z}
\DeclareUnicodeCharacter{0397}{H}
\DeclareUnicodeCharacter{0398}{\ensuremath{\Theta}}
\DeclareUnicodeCharacter{0399}{I}
\DeclareUnicodeCharacter{039A}{K}
\DeclareUnicodeCharacter{039B}{\ensuremath{\Lambda}}
\DeclareUnicodeCharacter{039C}{M}
\DeclareUnicodeCharacter{039D}{N}
\DeclareUnicodeCharacter{039E}{\ensuremath{\Xi}}
\DeclareUnicodeCharacter{039F}{O}
\DeclareUnicodeCharacter{03A0}{\ensuremath{\Pi}}
\DeclareUnicodeCharacter{03A1}{P}
\DeclareUnicodeCharacter{03A3}{\ensuremath{\Sigma}}
\DeclareUnicodeCharacter{03A4}{T}
\DeclareUnicodeCharacter{03A5}{Y}
\DeclareUnicodeCharacter{03A6}{\ensuremath{\Phi}}
\DeclareUnicodeCharacter{03A7}{X}
\DeclareUnicodeCharacter{03A8}{\ensuremath{\Psi}}
\DeclareUnicodeCharacter{03A9}{\ensuremath{\Omega}}
\DeclareUnicodeCharacter{03B1}{\ensuremath{\alpha}}
\DeclareUnicodeCharacter{03B2}{\ensuremath{\beta}}
\DeclareUnicodeCharacter{03B3}{\ensuremath{\gamma}}
\DeclareUnicodeCharacter{03B4}{\ensuremath{\delta}}
\DeclareUnicodeCharacter{03B5}{\ensuremath{\epsilon}}
\DeclareUnicodeCharacter{03B6}{\ensuremath{\zeta}}
\DeclareUnicodeCharacter{03B7}{\ensuremath{\eta}}
\DeclareUnicodeCharacter{03B8}{\ensuremath{\theta}}
\DeclareUnicodeCharacter{03B9}{\ensuremath{\iota}}
\DeclareUnicodeCharacter{03BA}{\ensuremath{\kappa}}
\DeclareUnicodeCharacter{03BB}{\ensuremath{\lambda}}
\DeclareUnicodeCharacter{03BC}{\ensuremath{\mu}}
\DeclareUnicodeCharacter{03BD}{\ensuremath{\nu}}
\DeclareUnicodeCharacter{03BE}{\ensuremath{\xi}}
\DeclareUnicodeCharacter{03BF}{\ensuremath{o}}
\DeclareUnicodeCharacter{03C0}{\ensuremath{\pi}}
\DeclareUnicodeCharacter{03C1}{\ensuremath{\rho}}
\DeclareUnicodeCharacter{03C2}{\ensuremath{\varsigma}}
\DeclareUnicodeCharacter{03C3}{\ensuremath{\sigma}}
\DeclareUnicodeCharacter{03C4}{\ensuremath{\tau}}
\DeclareUnicodeCharacter{03C5}{\ensuremath{\upsilon}}
\DeclareUnicodeCharacter{03C6}{\ensuremath{\phi}}
\DeclareUnicodeCharacter{03C7}{\ensuremath{\chi}}
\DeclareUnicodeCharacter{03C8}{\ensuremath{\psi}}
\DeclareUnicodeCharacter{03C9}{\ensuremath{\omega}}
\DeclareUnicodeCharacter{03D0}{\ensuremath{\beta}}
\DeclareUnicodeCharacter{03D1}{\ensuremath{\vartheta}}
\DeclareUnicodeCharacter{03D5}{\ensuremath{\varphi}}
\DeclareUnicodeCharacter{03D6}{\ensuremath{\varpi}}
\DeclareUnicodeCharacter{03F0}{\ensuremath{\kappa}}
\DeclareUnicodeCharacter{03F1}{\ensuremath{\varrho}}
\DeclareUnicodeCharacter{03F5}{\ensuremath{\varepsilon}}
\DeclareUnicodeCharacter{00D7}{\ensuremath{\times}}
\DeclareUnicodeCharacter{00B7}{\ensuremath{\cdot}}
\DeclareUnicodeCharacter{2212}{\ensuremath{-}}
\DeclareUnicodeCharacter{221E}{\ensuremath{\infty}}
\DeclareUnicodeCharacter{2192}{\ensuremath{\to}}
\DeclareUnicodeCharacter{2264}{\ensuremath{\leq}}
\DeclareUnicodeCharacter{2265}{\ensuremath{\geq}}
\DeclareUnicodeCharacter{2260}{\ensuremath{\neq}}
\DeclareUnicodeCharacter{2013}{--}
\DeclareUnicodeCharacter{2014}{---}
\DeclareUnicodeCharacter{2018}{\textquoteleft}
\DeclareUnicodeCharacter{2019}{\textquoteright}
\DeclareUnicodeCharacter{201C}{\textquotedblleft}
\DeclareUnicodeCharacter{201D}{\textquotedblright}
\begin{document}
\begin{center}{\Large\bfseries Matrices from Rules}\\[0.35em]{\large Exact construction, symbolic algebra, and limits}\end{center}
\vspace{0.25em}\hrule\vspace{1em}
A matrix rule evaluates one expression over one-based row and column indices.

\section*{An exact Hilbert matrix}
Define the \(n\times n\) Hilbert matrix by
\[H(n)=\left(\frac{1}{i+j-1}\right)_{i,j}.\]

At \(n=3\), the rule expands to

\[H\left(3\right) = \begin{pmatrix}1 & \frac{1}{2} & \frac{1}{3} \\ \frac{1}{2} & \frac{1}{3} & \frac{1}{4} \\ \frac{1}{3} & \frac{1}{4} & \frac{1}{5}\end{pmatrix}\]

Its determinant stays exact:

\[\begin{aligned}H\left(3\right) &= \begin{pmatrix}1 & \frac{1}{2} & \frac{1}{3} \\ \frac{1}{2} & \frac{1}{3} & \frac{1}{4} \\ \frac{1}{3} & \frac{1}{4} & \frac{1}{5}\end{pmatrix} \\
&\sim \begin{pmatrix}1 & \frac{1}{2} & \frac{1}{3} \\ 0 & \frac{\frac{1}{4}}{3} & \frac{\frac{1}{3}}{4} \\ \frac{1}{3} & \frac{1}{4} & \frac{1}{5}\end{pmatrix} && \left(R_{2}\leftarrow R_{2}-\left(\frac{1}{2}\right)R_{1}\right) \\
&\sim \begin{pmatrix}1 & \frac{1}{2} & \frac{1}{3} \\ 0 & \frac{\frac{1}{4}}{3} & \frac{\frac{1}{3}}{4} \\ 0 & \frac{\frac{1}{3}}{4} & \frac{\frac{4}{9}}{5}\end{pmatrix} && \left(R_{3}\leftarrow R_{3}-\left(\frac{1}{3}\right)R_{1}\right) \\
&\sim \begin{pmatrix}1 & \frac{1}{2} & \frac{1}{3} \\ 0 & \frac{\frac{1}{4}}{3} & \frac{\frac{1}{3}}{4} \\ 0 & 0 & \frac{\frac{1}{36}}{5}\end{pmatrix} && \left(R_{3}\leftarrow R_{3}-\left(\frac{\frac{\frac{1}{3}}{4}}{\frac{\frac{1}{4}}{3}}\right)R_{2}\right) \\
\det\left(H\left(3\right)\right) &= 1\cdot \frac{\frac{1}{4}}{3}\cdot \frac{\frac{1}{36}}{5} \\
&= \frac{1}{2160}\end{aligned}\]

Gauss--Jordan elimination gives an integer inverse:

\[\begin{aligned}\left[H\left(3\right)\mid I\right] &= \begin{pmatrix}1 & \frac{1}{2} & \frac{1}{3} & 1 & 0 & 0 \\ \frac{1}{2} & \frac{1}{3} & \frac{1}{4} & 0 & 1 & 0 \\ \frac{1}{3} & \frac{1}{4} & \frac{1}{5} & 0 & 0 & 1\end{pmatrix} \\
&\sim \begin{pmatrix}1 & \frac{1}{2} & \frac{1}{3} & 1 & 0 & 0 \\ 0 & \frac{\frac{1}{4}}{3} & \frac{\frac{1}{3}}{4} & -\left(\frac{1}{2}\right) & 1 & 0 \\ \frac{1}{3} & \frac{1}{4} & \frac{1}{5} & 0 & 0 & 1\end{pmatrix} && \left(R_{2}\leftarrow R_{2}-\left(\frac{1}{2}\right)R_{1}\right) \\
&\sim \begin{pmatrix}1 & \frac{1}{2} & \frac{1}{3} & 1 & 0 & 0 \\ 0 & \frac{\frac{1}{4}}{3} & \frac{\frac{1}{3}}{4} & -\left(\frac{1}{2}\right) & 1 & 0 \\ 0 & \frac{\frac{1}{3}}{4} & \frac{\frac{4}{9}}{5} & -\left(\frac{1}{3}\right) & 0 & 1\end{pmatrix} && \left(R_{3}\leftarrow R_{3}-\left(\frac{1}{3}\right)R_{1}\right) \\
&\sim \begin{pmatrix}1 & \frac{1}{2} & \frac{1}{3} & 1 & 0 & 0 \\ 0 & 1 & \frac{\frac{\frac{1}{3}}{4}}{\frac{\frac{1}{4}}{3}} & \frac{-\left(\frac{1}{2}\right)}{\frac{\frac{1}{4}}{3}} & \frac{1}{\frac{\frac{1}{4}}{3}} & 0 \\ 0 & \frac{\frac{1}{3}}{4} & \frac{\frac{4}{9}}{5} & -\left(\frac{1}{3}\right) & 0 & 1\end{pmatrix} && \left(R_{2}\leftarrow \frac{1}{\frac{\frac{1}{4}}{3}}R_{2}\right) \\
&\sim \begin{pmatrix}1 & 0 & \frac{-\frac{1}{2}}{3} & 4 & -6 & 0 \\ 0 & 1 & \frac{\frac{\frac{1}{3}}{4}}{\frac{\frac{1}{4}}{3}} & \frac{-\left(\frac{1}{2}\right)}{\frac{\frac{1}{4}}{3}} & \frac{1}{\frac{\frac{1}{4}}{3}} & 0 \\ 0 & \frac{\frac{1}{3}}{4} & \frac{\frac{4}{9}}{5} & -\left(\frac{1}{3}\right) & 0 & 1\end{pmatrix} && \left(R_{1}\leftarrow R_{1}-\left(\frac{1}{2}\right)R_{2}\right) \\
&\sim \begin{pmatrix}1 & 0 & \frac{-\frac{1}{2}}{3} & 4 & -6 & 0 \\ 0 & 1 & \frac{\frac{\frac{1}{3}}{4}}{\frac{\frac{1}{4}}{3}} & \frac{-\left(\frac{1}{2}\right)}{\frac{\frac{1}{4}}{3}} & \frac{1}{\frac{\frac{1}{4}}{3}} & 0 \\ 0 & 0 & \frac{\frac{1}{36}}{5} & \frac{1}{6} & -1 & 1\end{pmatrix} && \left(R_{3}\leftarrow R_{3}-\left(\frac{\frac{1}{3}}{4}\right)R_{2}\right) \\
&\sim \begin{pmatrix}1 & 0 & \frac{-\frac{1}{2}}{3} & 4 & -6 & 0 \\ 0 & 1 & \frac{\frac{\frac{1}{3}}{4}}{\frac{\frac{1}{4}}{3}} & \frac{-\left(\frac{1}{2}\right)}{\frac{\frac{1}{4}}{3}} & \frac{1}{\frac{\frac{1}{4}}{3}} & 0 \\ 0 & 0 & 1 & \frac{\frac{1}{6}}{\frac{\frac{1}{36}}{5}} & \frac{-1}{\frac{\frac{1}{36}}{5}} & \frac{1}{\frac{\frac{1}{36}}{5}}\end{pmatrix} && \left(R_{3}\leftarrow \frac{1}{\frac{\frac{1}{36}}{5}}R_{3}\right) \\
&\sim \begin{pmatrix}1 & 0 & 0 & 9 & -36 & 30 \\ 0 & 1 & \frac{\frac{\frac{1}{3}}{4}}{\frac{\frac{1}{4}}{3}} & \frac{-\left(\frac{1}{2}\right)}{\frac{\frac{1}{4}}{3}} & \frac{1}{\frac{\frac{1}{4}}{3}} & 0 \\ 0 & 0 & 1 & \frac{\frac{1}{6}}{\frac{\frac{1}{36}}{5}} & \frac{-1}{\frac{\frac{1}{36}}{5}} & \frac{1}{\frac{\frac{1}{36}}{5}}\end{pmatrix} && \left(R_{1}\leftarrow R_{1}-\left(\frac{-\frac{1}{2}}{3}\right)R_{3}\right) \\
&\sim \begin{pmatrix}1 & 0 & 0 & 9 & -36 & 30 \\ 0 & 1 & 0 & \frac{-3}{\frac{\frac{1}{4}}{3}} & \frac{16}{\frac{\frac{1}{4}}{3}} & -180 \\ 0 & 0 & 1 & \frac{\frac{1}{6}}{\frac{\frac{1}{36}}{5}} & \frac{-1}{\frac{\frac{1}{36}}{5}} & \frac{1}{\frac{\frac{1}{36}}{5}}\end{pmatrix} && \left(R_{2}\leftarrow R_{2}-\left(\frac{\frac{\frac{1}{3}}{4}}{\frac{\frac{1}{4}}{3}}\right)R_{3}\right) \\
H\left(3\right)^{-1} &= \begin{pmatrix}9 & -36 & 30 \\ -36 & 192 & -180 \\ 30 & -180 & 180\end{pmatrix}\end{aligned}\]

\section*{Symbolic entries and matrix-valued functions}
An explicit dimension lets the function parameter remain inside the entries:
\[S(t)=\left(t^{i+j-2}\right)_{i,j=1}^{2}.\]

\[S\left(t\right) = \begin{pmatrix}1 & t \\ t & t^{2}\end{pmatrix}\]

The rows are symbolically dependent, so the determinant vanishes:

\[\begin{aligned}S\left(t\right) &= \begin{pmatrix}1 & t \\ t & t^{2}\end{pmatrix} \\
&\sim \begin{pmatrix}1 & t \\ 0 & 0\end{pmatrix} && \left(R_{2}\leftarrow R_{2}-\left(t\right)R_{1}\right) \\
\det\left(S\left(t\right)\right) &= 0\end{aligned}\]

Matrix limits are taken entry by entry:

\[\lim_{t\to 0} S\left(t\right) = \begin{pmatrix}1 & 0 \\ 0 & 0\end{pmatrix}\]

A rectangular rule supplies both dimensions explicitly:

\[\left(10i + j\right)_{i,j}^{2\times3} = \begin{pmatrix}11 & 12 & 13 \\ 21 & 22 & 23\end{pmatrix}\]

\section*{Limits through linear algebra}
For the averaging matrix \(P\), powers converge to projection onto the steady state.

\[\lim_{n\to \infty} P^{n} = \begin{pmatrix}\frac{1}{2} & \frac{1}{2} \\ \frac{1}{2} & \frac{1}{2}\end{pmatrix}\]

Continuous linear-algebra operations consume the same symbolic sequence:

\begin{align*}\lim_{n\to \infty} \operatorname{tr}\left(P^{n}\right) = 1 \\
\lim_{n\to \infty} \det\left(P^{n}\right) = 0\end{align*}

The source definition, exact calculations, and displayed matrices above are all evaluated by SCAS while the document is built.
\end{document}
